% planisphere_ru.tex
%
% The LaTeX code in this file brings together into a single document the
% various components of the model planisphere.
%
% Copyright (C) 2014-2024 Dominic Ford <https://dcford.org.uk/>
%
% This code is free software; you can redistribute it and/or modify it under
% the terms of the GNU General Public License as published by the Free Software
% Foundation; either version 2 of the License, or (at your option) any later
% version.
%
% You should have received a copy of the GNU General Public License along with
% this file; if not, write to the Free Software Foundation, Inc., 51 Franklin
% Street, Fifth Floor, Boston, MA  02110-1301, USA

% ----------------------------------------------------------------------------

\documentclass[a4paper,onecolumn,10pt]{article}
\usepackage[dvips]{graphicx}
\usepackage{fancyhdr,url}
\usepackage[utf8]{inputenc}
\usepackage[T2A]{fontenc}
\usepackage[russian]{babel}
\usepackage{parskip}
\usepackage[pdftitle={Сделайте свою планисферу}, pdfauthor={Dominic Ford}, pdfsubject={Сделайте свою планисферу}, pdfkeywords={Сделайте свою планисферу}, colorlinks=true, linkcolor=blue, citecolor=blue, filecolor=blue, urlcolor=blue]{hyperref}
\renewcommand{\familydefault}{\sfdefault}
\pagestyle{fancy}

\lhead{\it СДЕЛАЙТЕ СВОЮ ПЛАНИСФЕРУ}
\chead{}
\rhead{\thepage}
\lfoot{}\rfoot{}
\cfoot{\bf\footnotesize\copyright\ 2014--2024 Dominic Ford. Распространяется по GNU General Public License, версия 3. Документ загружен с \url{https://in-the-sky.org/planisphere/}}

\fancypagestyle{plain}{%
\fancyhf{}
\renewcommand{\headrulewidth}{0pt}
\renewcommand{\footrulewidth}{0pt}}

\title{Сделайте свою планисферу}
\author{Dominic Ford}
\date{2014--2024}

\begin{document}
\maketitle
\setcounter{footnote}{1}

Планисфера --- это простое ручное устройство, показывающее карту звезд,
видимых на ночном небе в выбранный момент времени. Поворачивая диск,
можно увидеть, как звезды перемещаются по небу в течение ночи и как
разные созвездия становятся видимыми в разное время года.

Здесь представлен набор шаблонов, который можно скачать и распечатать,
чтобы собрать собственную планисферу из бумаги или тонкого картона.

Конструкция планисферы зависит от географической широты, где она будет
использоваться, поскольку в разных местах видны разные участки неба.
Готовые наборы для широкого диапазона широт доступны по адресу

\url{https://in-the-sky.org/planisphere/}

Планисфера в этом документе рассчитана на широту
\input{tmp/lat}.

\section*{Что потребуется}

\begin{itemize}
\item Два листа A4, лучше тонкий картон.
\item Ножницы.
\item Канцелярская кнопка-брадс.
\item Необязательно: лист прозрачного пластика (например, пленка для проекторов).
\item Необязательно: немного клея.
\end{itemize}

\section*{Инструкция по сборке}

{\bf Шаг 1} --- Планисферы отличаются в зависимости от места использования.
Данный набор рассчитан на широты вблизи \input{tmp/lat}. Если вы живете
на другой широте, скачайте подходящий набор с сайта:

\url{https://in-the-sky.org/planisphere/}

{\bf Шаг 2} --- Распечатайте страницы в конце PDF со звездным диском и
корпусом планисферы на двух отдельных листах бумаги, а лучше --- на тонком
картоне.

{\bf Шаг 3} --- Аккуратно вырежьте звездный диск и корпус планисферы.
Также вырежьте серую затененную область на корпусе и, при необходимости,
лист с сеткой линий, напечатанный на прозрачном пластике. Если используете
картон, можно слегка продавить линию сгиба по пунктиру.

{\bf Шаг 4} --- В центре звездного диска и внизу корпуса есть маленькие
круглые метки. Сделайте в этих местах небольшие отверстия (примерно 2 мм).
Удобно использовать дырокол для бумаги, а при его отсутствии --- острие
циркуля.

{\bf Шаг 5} --- Вставьте брадс в центр звездного диска со стороны печати.
Затем наденьте на тот же брадс корпус планисферы так, чтобы печатная
сторона корпуса была обращена к обратной стороне брадса. Разведите ножки
брадса, чтобы закрепить обе детали.

{\bf Шаг 6 (необязательно)} --- Если последняя страница PDF была распечатана
на прозрачной пленке, наклейте эту сетку поверх вырезанного окна на корпусе
планисферы.

{\bf Шаг 7} --- Согните корпус планисферы по пунктирной линии так, чтобы
лицевая сторона звездного диска была видна через вырезанное окно.

{\bf Готово! Ваша планисфера готова к использованию.}

\section*{Как пользоваться планисферой}

Поверните звездный диск так, чтобы сегодняшняя дата на его краю совпала
с текущим временем. В смотровом окне будут показаны созвездия, видимые
сейчас на небе.

Выйдите на улицу и повернитесь к северу (в северном полушарии) или к югу
(в южном полушарии). Поднимите планисферу к небу: звезды в нижней части
окна должны совпадать с теми, что вы видите перед собой.

Повернитесь к востоку или западу и поверните планисферу так, чтобы надпись
"Восток" или "Запад" оказалась в нижней части окна. Нижняя часть окна снова
должна совпадать с видимыми впереди звездами.

Если вы распечатали сетку высот и азимутов на прозрачной пленке, по ней
можно оценивать высоту объектов над горизонтом и направление на них. Круги
соответствуют высотам 10, 20, 30, ..., 80 градусов. Для ориентира: 10
градусов примерно соответствуют ширине ладони на вытянутой руке. Изогнутые
линии показывают вертикали для направлений S, SSE, SE, ESE, E и т. д.

\section*{Кастомизация}

Этот набор планисферы создан с помощью Python-скриптов и библиотеки
pycairo. Если хотите настроить планисферу под себя, вы можете скачать
исходный код с GitHub и изменить его с указанием авторства:

\url{https://github.com/dcf21/planisphere}

\section*{Лицензия}

Как и весь контент на {\tt In-The-Sky.org}, этот набор защищен авторским
правом \copyright\ Dominic Ford. При этом материалы доступны для любителей
астрономии по всему миру, и вы можете изменять и/или распространять их при
условиях: (1) любой материал с текстом об авторских правах {\bf должен}
содержать этот текст {\bf без изменений}; (2) вы {\bf должны} указывать
Dominic Ford как исходного автора и правообладателя; (3) вы {\bf не можете}
получать прибыль от распространения материалов сайта, {\bf если только} вы
не являетесь официально зарегистрированной благотворительной организацией,
продвигающей астрономическую науку, {\bf или} не получили письменное
разрешение автора.

\newpage

\centerline{\includegraphics{tmp/starwheel}}

\vspace{1cm}
Центральный звездный диск планисферы, который должен располагаться внутри сложенного корпуса.

\newpage
\thispagestyle{empty}
\vspace*{-3.0cm}
\centerline{\includegraphics{tmp/holder}}
\newpage

\centerline{\includegraphics{tmp/altaz}}

\vspace{1cm}
Эту сетку можно дополнительно распечатать на прозрачной пленке и наклеить
в вырезанное окно корпуса планисферы для оценки высоты объектов на небе и
направления на них.

\end{document}
